(1) 一条中心位于 x0 的宽 a 狭缝远场振幅正比
\[
A_1(\theta)=E_0\int_{-a/2}^{a/2}e^{ikx\sin\theta}\,dx
=E_0a\frac{\sin\beta}{\beta},
\quad \beta=\frac{\pi a\sin\theta}{\lambda}.
\]
其强度包络为 \(\mathrm{sinc}^2\beta\)。

(2) 两缝中心相距 d，复振幅之和多出 \(2\cos\delta\) 因子，其中 \(\delta=\pi d\sin\theta/\lambda\)。归一化总强度
\[
I(\theta)=4I_1(0)\left(\frac{\sin\beta}{\beta}\right)^2\cos^2\delta.
\]
中心 \(\theta =0\) 处强度是一个狭缝峰值的4倍，因两缝场同相叠加。

(3) 双缝相长位置 \(d\sin\theta=m\lambda,\ m\in\mathbb Z\)。单缝衍射零点位置 \(a\sin\theta=j\lambda\)，其中 j 为非零整数。两个条件分别来自干涉因子和单缝振幅，不能互相代替。

(4) 若某干涉主极大恰满足单缝零点，则该级消失。联立给 \(m a/d=j\)，即 \(m=j d/a\) 必须是整数。特别当 d/a 为整数 p，所有非零的第 p、2p、3p……级干涉峰缺失。此结论假设两缝相同且完全相干，真实有限孔径与对准误差会使零点不完全消光。
